% ============================================================
%  SEL RTAC — COURSE SYLLABUS
% ============================================================
\documentclass[10pt]{article}
\usepackage[T1]{fontenc}
\usepackage[utf8]{inputenc}
\usepackage{lmodern}
\usepackage{microtype}
\usepackage{palatino}
\usepackage[scaled=0.82]{beramono}
\usepackage[top=0.55in, bottom=0.50in, left=0.65in, right=0.65in]{geometry}
\usepackage[dvipsnames,svgnames,x11names]{xcolor}
\usepackage{tikz}
\usepackage{enumitem}
\usepackage[most]{tcolorbox}
\usepackage[colorlinks=true, linkcolor=rgreen, urlcolor=rcyan,
  pdftitle={SEL RTAC --- Course Syllabus}]{hyperref}
\usepackage{parskip}

\definecolor{rgreen}{HTML}{15803D}
\definecolor{rcyan}{HTML}{059669}
\definecolor{rnavy}{HTML}{052E16}
\definecolor{rorange}{HTML}{D97706}
\definecolor{raccent}{HTML}{F0FDF4}
\definecolor{rgray}{HTML}{6B7280}

\setlength{\parskip}{2pt}
\setlength{\parindent}{0pt}
\pagestyle{empty}

\begin{document}

\begin{tikzpicture}[remember picture, overlay]
  \fill[rnavy] (current page.north west) rectangle ([yshift=-1.0in]current page.north east);
  \node[anchor=west, text=white, font=\fontsize{22}{26}\bfseries\sffamily]
    at ([xshift=0.65in, yshift=-0.35in]current page.north west) {SEL RTAC};
  \node[anchor=west, text=rcyan!80, font=\fontsize{11}{13}\sffamily]
    at ([xshift=0.65in, yshift=-0.63in]current page.north west)
    {Industrial Control Systems · SCADA Hub};
  \draw[rorange, line width=2pt]
    ([xshift=0.65in, yshift=-0.83in]current page.north west) --
    ([xshift=3.5in,  yshift=-0.83in]current page.north west);
\end{tikzpicture}

\vspace{0.95in}

{\color{rnavy}\rule{\linewidth}{1.2pt}}
\smallskip
{\small\sffamily\color{rgray}\MakeUppercase{Course Topics}}
\smallskip

\setlist[enumerate]{noitemsep, topsep=2pt, leftmargin=*, label=\textcolor{rgreen}{\small\bfseries\arabic*.}}
\begin{enumerate}\small
  \item \textbf{RTAC Architecture:} SEL-3505/3530/3555 hardware variants, I/O module selection, communication ports, and physical commissioning steps.
  \item \textbf{ACSELERATOR Architect:} Project database structure, tags, aliases, protocol point mapping, revision control, and first-live-data workflow.
  \item \textbf{DNP3 Configuration:} Outstation and master roles, point database mapping, event class assignment, unsolicited reporting, and SCADA latency parameters.
  \item \textbf{Modbus Gateway:} TCP server configuration, DNP3-to-Modbus point mapping, register layout, data type conversion, and translation failure modes.
  \item \textbf{IEC 61850 Integration:} GOOSE trip signals, Sampled Values, MMS operator control, and internal tag database mapping for each service type.
  \item \textbf{Local Logic:} IEC 61131-3 programs, scan cycle behavior, event-driven automation, and design principles for WAN-independent operation.
  \item \textbf{Security and Access Control:} Role-based access, serial port security, firmware update procedures, and NERC CIP implications per configuration decision.
  \item \textbf{Field Scenarios:} Three case studies — automatic tap changer response, feeder protection with downstream reclosers, and sub-200ms load shedding.
\end{enumerate}

\medskip
{\color{rnavy}\rule{\linewidth}{0.4pt}}
\smallskip
{\small\sffamily\color{rgray}\MakeUppercase{Learning Outcomes}}
\smallskip

\setlist[itemize]{noitemsep, topsep=2pt, leftmargin=*, label=\textcolor{rgreen}{\small$\bullet$}}
\begin{itemize}\small
  \item Commission an SEL RTAC from hardware selection through first live DNP3 tag.
  \item Map DNP3, Modbus, and IEC 61850 protocol points to the RTAC internal tag database.
  \item Write IEC 61131-3 local logic programs that operate correctly without WAN connectivity.
  \item Configure role-based access control and document each decision against NERC CIP requirements.
  \item Analyze and implement field automation scenarios including tap changer and load shedding sequences.
\end{itemize}

\vfill
{\color{rgreen}\rule{\linewidth}{0.6pt}}
\smallskip

\begin{minipage}[t]{0.48\textwidth}
\begin{tcolorbox}[
  enhanced, colback=raccent, colframe=rgreen, arc=4pt,
  fonttitle=\small\bfseries\sffamily\color{white},
  title=Certification Relevance,
  boxed title style={colback=rgreen},
  attach boxed title to top left={yshift=-2mm, xshift=5mm},
  top=6pt, bottom=4pt, left=5pt, right=5pt
]
\setlist[itemize]{noitemsep, topsep=0pt, leftmargin=*, label=\textcolor{rgreen}{$\bullet$}}
\begin{itemize}\small
  \item \textbf{SEL Certified Engineer} — SEL offers application engineer certification through their training center; RTAC configuration is a core track.
  \item \textbf{GICSP (GIAC)} — Substation automation, IEC 61850 security, and protocol gateway configuration appear in ICS security exam domains.
  \item \textbf{ISA CCST} — Protocol gateway integration and field device configuration tested at CCST Levels II and III.
\end{itemize}
\end{tcolorbox}
\end{minipage}
\hfill
\begin{minipage}[t]{0.48\textwidth}
\begin{tcolorbox}[
  enhanced, colback=rorange!8, colframe=rorange, arc=4pt,
  fonttitle=\small\bfseries\sffamily\color{white},
  title=Next Steps,
  boxed title style={colback=rorange},
  attach boxed title to top left={yshift=-2mm, xshift=5mm},
  top=6pt, bottom=4pt, left=5pt, right=5pt
]
\begin{enumerate}[noitemsep, topsep=0pt, leftmargin=*, label=\small\textcolor{rorange}{\arabic*.}]
  \item \small Take DNP3 and IEC 61131-3 before RTAC. Protocol knowledge and programming skills are both prerequisites for field configuration work.
  \item Download ACSELERATOR Architect (free from selinc.com) and configure a lab project against a virtual RTAC.
  \item Attend SEL application engineering training in Pullman, WA or at a regional event. The in-person lab is the fastest path to real proficiency.
  \item GICSP pairs well with RTAC expertise for engineers working in power utilities and substations.
\end{enumerate}
\end{tcolorbox}
\end{minipage}

\begin{tikzpicture}[remember picture, overlay]
  \fill[rnavy!90] (current page.south west) rectangle ([yshift=0.30in]current page.south east);
  \node[anchor=west, text=white!70, font=\footnotesize\sffamily]
    at ([xshift=0.65in, yshift=0.15in]current page.south west)
    {SCADA Hub · Industrial Control Systems · scada-hub.vercel.app};
  \node[anchor=east, text=rcyan, font=\footnotesize\bfseries\sffamily]
    at ([xshift=-0.65in, yshift=0.15in]current page.south east)
    {Course 6 of 8};
\end{tikzpicture}

\end{document}
